%Version 3 December 2023
% See section 11 of the User Manual for version history
%
%%%%%%%%%%%%%%%%%%%%%%%%%%%%%%%%%%%%%%%%%%%%%%%%%%%%%%%%%%%%%%%%%%%%%%
%%                                                                 %%
%% Please do not use \input{...} to include other tex files.       %%
%% Submit your LaTeX manuscript as one .tex document.              %%
%%                                                                 %%
%% All additional figures and files should be attached             %%
%% separately and not embedded in the \TeX\ document itself.       %%
%%                                                                 %%
%%%%%%%%%%%%%%%%%%%%%%%%%%%%%%%%%%%%%%%%%%%%%%%%%%%%%%%%%%%%%%%%%%%%%

%%\documentclass[referee,sn-basic]{sn-jnl}% referee option is meant for double line spacing

%%=======================================================%%
%% to print line numbers in the margin use lineno option %%
%%=======================================================%%

%%\documentclass[lineno,sn-basic]{sn-jnl}% Basic Springer Nature Reference Style/Chemistry Reference Style

%%======================================================%%
%% to compile with pdflatex/xelatex use pdflatex option %%
%%======================================================%%

%%\documentclass[pdflatex,sn-basic]{sn-jnl}% Basic Springer Nature Reference Style/Chemistry Reference Style


%%Note: the following reference styles support Namedate and Numbered referencing. By default the style follows the most common style. To switch between the options you can add or remove “Numbered” in the optional parenthesis. 
%%The option is available for: sn-basic.bst, sn-vancouver.bst, sn-chicago.bst%  
 
%%\documentclass[pdflatex,sn-nature]{sn-jnl}% Style for submissions to Nature Portfolio journals
%%\documentclass[pdflatex,sn-basic]{sn-jnl}% Basic Springer Nature Reference Style/Chemistry Reference Style
\documentclass[pdflatex,sn-mathphys-num]{sn-jnl}% Math and Physical Sciences Numbered Reference Style 
%%\documentclass[pdflatex,sn-mathphys-ay]{sn-jnl}% Math and Physical Sciences Author Year Reference Style
%%\documentclass[pdflatex,sn-aps]{sn-jnl}% American Physical Society (APS) Reference Style
%%\documentclass[pdflatex,sn-vancouver,Numbered]{sn-jnl}% Vancouver Reference Style
%%\documentclass[pdflatex,sn-apa]{sn-jnl}% APA Reference Style 
%%\documentclass[pdflatex,sn-chicago]{sn-jnl}% Chicago-based Humanities Reference Style

%%%% Standard Packages
%%<additional latex packages if required can be included here>

\usepackage{graphicx}%
\usepackage{multirow}%
\usepackage{amsmath,amssymb,amsfonts}%
\usepackage{amsthm}%
\usepackage{mathrsfs}%
\usepackage[title]{appendix}%
\usepackage{xcolor}%
\usepackage{textcomp}%
\usepackage{manyfoot}%
\usepackage{booktabs}%
\usepackage{algorithm}%
\usepackage{algorithmicx}%
\usepackage{algpseudocode}%
\usepackage{listings}%
%%%%
\usepackage{caption}


% \usepackage{verbatim} % For comment environment
% \renewenvironment{figure}[1][]{%
%   \expandafter\comment%
% }{%
%   \expandafter\endcomment%
% }
% \renewenvironment{table}[1][]{%
%   \expandafter\comment%
% }{%
%   \expandafter\endcomment%
% }


\newcommand*{\transpose}{^{\top}}

%%%%%=============================================================================%%%%
%%%%  Remarks: This template is provided to aid authors with the preparation
%%%%  of original research articles intended for submission to journals published 
%%%%  by Springer Nature. The guidance has been prepared in partnership with 
%%%%  production teams to conform to Springer Nature technical requirements. 
%%%%  Editorial and presentation requirements differ among journal portfolios and 
%%%%  research disciplines. You may find sections in this template are irrelevant 
%%%%  to your work and are empowered to omit any such section if allowed by the 
%%%%  journal you intend to submit to. The submission guidelines and policies 
%%%%  of the journal take precedence. A detailed User Manual is available in the 
%%%%  template package for technical guidance.
%%%%%=============================================================================%%%%

%% as per the requirement new theorem styles can be included as shown below
%\theoremstyle{thmstyleone}%
%\newtheorem{theorem}{Theorem}%  meant for continuous numbers
%%\newtheorem{theorem}{Theorem}[section]% meant for sectionwise numbers
%% optional argument [theorem] produces theorem numbering sequence instead of independent numbers for Proposition
%\newtheorem{proposition}[theorem]{Proposition}% 
%%\newtheorem{proposition}{Proposition}% to get separate numbers for theorem and proposition etc.

%\theoremstyle{thmstyletwo}%
%\newtheorem{example}{Example}%
\newtheorem*{remark}{Remark}

%\theoremstyle{thmstylethree}%
%\newtheorem{definition}{Definition}%

%\raggedbottom
%%\unnumbered% uncomment this for unnumbered level heads

\begin{document}

\title[Article Title]{Exact and approximate formulations for the close-enough TSP}

%%=============================================================%%
%% GivenName	-> \fnm{Joergen W.}
%% Particle	-> \spfx{van der} -> surname prefix
%% FamilyName	-> \sur{Ploeg}
%% Suffix	-> \sfx{IV}
%% \author*[1,2]{\fnm{Joergen W.} \spfx{van der} \sur{Ploeg} 
%%  \sfx{IV}}\email{iauthor@gmail.com}
%%=============================================================%%

\author*[1,2]{\fnm{First} \sur{Author}}\email{iauthor@gmail.com}

\author[2,3]{\fnm{Second} \sur{Author}}\email{iiauthor@gmail.com}
\equalcont{These authors contributed equally to this work.}

\author[1,2]{\fnm{Third} \sur{Author}}\email{iiiauthor@gmail.com}
\equalcont{These authors contributed equally to this work.}

\affil*[1]{\orgdiv{Department}, \orgname{Organization}, \orgaddress{\street{Street}, \city{City}, \postcode{100190}, \state{State}, \country{Country}}}

\affil[2]{\orgdiv{Department}, \orgname{Organization}, \orgaddress{\street{Street}, \city{City}, \postcode{10587}, \state{State}, \country{Country}}}

\affil[3]{\orgdiv{Department}, \orgname{Organization}, \orgaddress{\street{Street}, \city{City}, \postcode{610101}, \state{State}, \country{Country}}}

%%==================================%%
%% Sample for unstructured abstract %%
%%==================================%%

\abstract{The abstract serves both as a general introduction to the topic and as a brief, non-technical summary of the main results and their implications. Authors are advised to check the author instructions for the journal they are submitting to for word limits and if structural elements like subheadings, citations, or equations are permitted.}

%%================================%%
%% Sample for structured abstract %%
%%================================%%

% \abstract{\textbf{Purpose:} The abstract serves both as a general introduction to the topic and as a brief, non-technical summary of the main results and their implications. The abstract must not include subheadings (unless expressly permitted in the journal's Instructions to Authors), equations or citations. As a guide the abstract should not exceed 200 words. Most journals do not set a hard limit however authors are advised to check the author instructions for the journal they are submitting to.
% 
% \textbf{Methods:} The abstract serves both as a general introduction to the topic and as a brief, non-technical summary of the main results and their implications. The abstract must not include subheadings (unless expressly permitted in the journal's Instructions to Authors), equations or citations. As a guide the abstract should not exceed 200 words. Most journals do not set a hard limit however authors are advised to check the author instructions for the journal they are submitting to.
% 
% \textbf{Results:} The abstract serves both as a general introduction to the topic and as a brief, non-technical summary of the main results and their implications. The abstract must not include subheadings (unless expressly permitted in the journal's Instructions to Authors), equations or citations. As a guide the abstract should not exceed 200 words. Most journals do not set a hard limit however authors are advised to check the author instructions for the journal they are submitting to.
% 
% \textbf{Conclusion:} The abstract serves both as a general introduction to the topic and as a brief, non-technical summary of the main results and their implications. The abstract must not include subheadings (unless expressly permitted in the journal's Instructions to Authors), equations or citations. As a guide the abstract should not exceed 200 words. Most journals do not set a hard limit however authors are advised to check the author instructions for the journal they are submitting to.}

\keywords{keyword1, Keyword2, Keyword3, Keyword4}

%%\pacs[JEL Classification]{D8, H51}

%%\pacs[MSC Classification]{35A01, 65L10, 65L12, 65L20, 65L70}

\maketitle

\section{Introduction}\label{sec1}


The Traveling Salesman Problem (TSP) is a classic combinatorial optimization problem that has been the subject of extensive study for several decades. Given a set of locations and pairwise distances, the goal is to find the shortest possible tour that starts and ends at a designated location while visiting every other location exactly once. The TSP has several applications in a variety of fields, such as logistics, genome sequencing, and data clustering \cite{book:Applegate}.

This work handles the Close-Enough Traveling Salesman Problem (CETSP), a generalization of the TSP where, instead of it being necessary to reach the exact location, it is sufficient to visit a previously defined compact region, known as its neighborhood, that contains it. Although the problem can be defined for arbitrary compact regions, we assume that the neighborhoods are closed disks centered in the nodes, which is a common assumption in the CETSP literature \cite{behdani, branch}, and corresponds precisely to the type of neighborhoods that arise naturally in applications of the CETSP, which include distant meter readings via radio frequency identification technology \cite{genetic}, UAV routing in various contexts \cite{branch}, and robot monitoring of wireless sensor networks \cite{yuan}. 

Several generalizations of the CETSP have been proposed to capture additional practical requirements. For instance, \citet{generalized} propose the generalized CETSP, a variation of the problem where each location has several neighborhoods instead of a single one, with neighborhoods given by disks of different radii, and a prize associated with each region which is redeemed by traversing it. \citet{mcetsp} propose the multiple CETSP, where there are several vehicles, and the objective is to minimize the length of the longest tour assigned to one of the vehicles. \citet{walteros} study a variant of the CETSP called the pickup-and-delivery TSP with neighborhoods (PDTSPN), that combines traditional pickup and delivery requirements with the flexibility of the CETSP. In this context, they propose a cut-generation scheme for solving the problem, based on the generalized Benders decomposition \cite{gbd}.

While there are several heuristic approaches to solving the CETSP \cite{genetic, mennell, memetic}, there is a notorious scarcity of exact algorithms with optimality guarantees. \citet{mennell} introduced the first mathematical formulation of the problem, along with a heuristic approach based on finding non-empty neighborhood intersections, to subsequently solve a TSP over a selection of points from said intersections. Additionally, they introduced the \textit{overlap ratio} as a measure of instance difficulty, computed as the ratio between the mean radius and the largest side of the smallest rectangle that contains all the neighborhoods. \citet{gentilini2013travelling} presented a solution approach for this formulation based on a traditional spatial branch-and-bound algorithm for mixed-integer non-linear programming (MINLP), with specific improvements for the CETSP, in order to make the problem more tractable for MINLP solvers.

\citet{behdani} proposed a discretization scheme that generates lower and upper bounds on most test instances, and can be iteratively refined to make said bounds arbitrarily close. Starting from these results, \citet{carrabs} introduced a novel discretization scheme that improves upon the results of \cite{behdani}, generally providing better bounds in lesser times.

To the best of our knowledge, the only exact algorithm specific to this problem is a combinatorial \textit{branch-and-bound} \cite{branch}, where each node of the tree corresponds to a partial visit sequence, for which the visit points for each neighborhood are computed via second-order cone programming (SOCP). If the solution does not visit every neighborhood, the partial sequence is extended, creating new nodes of the branching tree. This approach provides good results for several benchmark instances, but, in the worst case, has to go through every single possible partial sequence. This algorithm was refined by \citet{bnb-upgrade}, proposing a number of improvements to it that lead to better results in most instances of the problem. \citet{bnb-efficient} further improved the algorithm by reusing information from parent nodes in the branching tree to speed up the solution of subproblems corresponding to the children nodes, while \citet{DECKEROVA2024125185} extended it to a more general setting.
 
In this paper, we present two SOCP-based formulations for the CETSP, along with approximated formulations based on mixed-integer linear programming (MILP) that provide arbitrarily close approximations to the feasible regions of the original formulations. Furthermore, we present decomposition schemes for both the SOCP and MILP formulations in order to make the problem more computationally manageable, and, in the case of the MILP models, to close the gap left by the approximation factor. We also provide examples that prove that the cuts proposed by \cite{walteros} are not always valid for the problem. Finally, we perform exhaustive computational experiments over a variety of instances to measure and compare the performance of the previously mentioned formulations.

The remainder of the paper is organized as follows. Section \ref{sec2} contains a formal definition of the CETSP, along with the notation that is used throughout the paper for the parameters of the problem, Section \ref{sec3} presents the SOCP-based exact formulations, Section \ref{sec4} introduces the linear approximation scheme and presents the resulting formulations, Section \ref{sec5} describes the two-stage decomposition scheme. In Section \ref{sec6}, we argue that the cuts proposed by \cite{walteros} are not actually valid, and we follow by presenting computational results in Section \ref{sec7}. Finally, conclusions are presented in Section \ref{sec8}.

\section{Problem description}\label{sec2}

The CETSP can be formally stated as follows. Let $N = \{1,\dots,n\}$ be a set of locations in a two-dimensional plane, with their corresponding coordinates $c_i = (c^x_i,c^y_i)$ for each $i \in N$. Each location $i \in N$ is covered by a closed disk centered in $c_i$ with radius $r_i$, which we call its neighborhood. In this setting, the CETSP seeks to find the shortest tour that visits each neighborhood, with no restrictions on visit order. We consider a neighborhood visited if the intersection between the tour and the neighborhood is not empty. Although the problem can be stated in terms of arbitrary norms, we consider Euclidean distances. We also assume $r_1 = 0$, so location $1$ acts as the depot.

\citet{behdani} proved that any optimal solution can be characterized by line segments between the neighborhoods, hence reducing the problem to finding a visit sequence, along with the endpoints of the corresponding line segments, in such a way that the tour over said points forms a Hamiltonian cycle of minimum length, and each point is contained in its corresponding neighborhood. Figure \ref{CETSP-example} illustrates an optimal CETSP tour for an instance with 5 nodes, exemplifying this property of optimal CETSP solutions.
\begin{figure}[H]
    \centering
    \includegraphics[width=0.8\textwidth]{cetsp_instance_plot.png}
    \caption{An optimal CETSP tour with $|N| = 5$.} \label{CETSP-example}
\end{figure}

\section{SOCP formulations}\label{sec3}

We now present our two SOCP formulations for the CETSP. Our first SOCP formulation is a natural generalization of the classic TSP formulation with a binary decision variable for each arc, which we call \textit{arc-based formulation} (\ref{abf}). In this formulation, variables $x_{ij} \in \{0,1\}$ indicate whether the tour goes from the neighborhood of location $i\in N$ to the neighborhood of location $j\in N$ or not. Variables $u_{ij}\geq 0$ are auxiliary variables utilized for subtour elimination, and variables $p_i \in \mathbb{R}^2$ represent the coordinates of the chosen point for the neighborhood corresponding to each location $i \in N$. Finally, variables $d_{ij}\geq 0$ represent the distance between the points chosen in each neighborhood, and variables $d_{ij}^a \geq 0$ are auxiliary variables that represent the distance of the line segments that compose the tour, taking the value 0 if the tour does not go from $i$ to $j$. 

The formulation is as follows:

\begin{subequations}
\begin{align}
\min &&& \sum_{i \in N}\sum_{j\in N} d_{ij}^a &  \tag{ABF} \label{abf}\\
\text{s.t.} &&& \sum_{j \in N} x_{ij} = 1 & \forall i \in N \label{tour-abf-1}\\
&&& \sum_{i \in N} x_{ij} = 1 & \forall j \in N \label{tour-abf-2}\\
&&& \sum_{j\in N} u_{ij} - \sum_{j \in N} u_{ji} = -1 & \forall i \in N\setminus \{1\} \label{sec-1}\\
&&& \sum_{j\in N}u_{1j} - \sum_{j\in N}u_{j1} = n - 1 & \label{sec-2}\\	
&&& u_{ij} \leq (n-1)x_{ij} &\forall (i,j) \in N\times N \label{sec-3}\\
&&& u_{ij} \geq 0 & \forall (i,j) \in N \times N \label{sec-4}\\
&&& ||p_i - c_i|| \leq r_i &\forall i \in N \label{neigh-abf}\\
&&& ||p_i - p_j|| \leq d_{ij} &\forall (i,j) \in N\times N \label{dist-abf}\\
&&& d_{ij} - M_{ij}(1 - x_{ij}) \leq d_{ij}^a &\forall (i,j) \in N\times N \label{big-m-abf}\\
&&& x_{ij} \in \{0,1\} &\forall (i,j) \in N\times N \label{x-abf}\\
&&& d_{ij}, d_{ij}^a \geq 0 &\forall (i,j) \in N\times N \label{d-abf}\\
&&& p_i \in \mathbb{R}^2 &\forall i \in N \label{p-abf}
\end{align}
\end{subequations}

The objective function of \ref{abf} minimizes the total distance of the tour. Constraints \eqref{tour-abf-1} and \eqref{tour-abf-2} are degree constraints, constraints \eqref{sec-1} to \eqref{sec-4} are subtour elimination constraints \cite{SEC}, constraint \eqref{neigh-abf} ensures that the points picked are contained in the neighborhoods, constraint \eqref{dist-abf} defines the distance between representative points, and constraint \eqref{big-m-abf} activates the variables $d_{ij}^a$, causing them to be equal to $x_{ij}d_{ij}$. Here, the parameter $M_{ij}$ corresponds to the maximum possible distance between neighborhoods $i$ and $j$, causing these constraints to be trivially satisfied if $x_{ij} = 0$. Lastly, constraints \eqref{sec-4}, \eqref{x-abf}-\eqref{p-abf} define the nature of the variables.

A potential problem with \ref{abf} is that it contains $\mathcal{O}(n^2)$ quadratic constraints, opening the possibility of the complexity rendering the problem intractable as $n$ increases. In addition, it relies on big-M constraints, which are known to have poor relaxations and to lead to numerical issues \cite{article:bigM}. This motivates our second formulation, based on determining which neighborhood gets visited at each of the $n$ stages of the tour. We refer to this formulation as \textit{sequence-based formulation} (\hyperref[sbf]{SBF}). In this formulation, variables $z_{ik} \in \{0,1\}$ indicate whether location $i \in N$ is visited in stage $k \in N$, $q_k \in \mathbb{R}^2$ represents the coordinates of the point visited in stage $k$, and $t_k \geq 0$ indicates the distance traveled in stage $k$.

The formulation is as follows:

\begin{subequations}
\begin{align}
\min &&& \sum_{k \in N} t_k \tag{SBF} \label{sbf}\\
\text{s.a.} &&& \sum_{k \in N} z_{ik} = 1 &\forall i \in N \label{tour-sbf-1}\\
&&& \sum_{i \in N} z_{ik} = 1 &\forall k \in N \label{tour-sbf-2}\\
&&& z_{11} = 1 \label{anti-sym}\\
&&& ||q_k - q_{k+1}|| \leq t_k &\forall k \in N \label{dist-sbf}\\
&&& \left|\left|\sum_{i \in N} {z}_{ik}  c_i - q_k\right|\right| \leq \sum_{i\in N} {z}_{ik}  r_i &\forall k \in N \label{neigh-sbf}\\
&&& z_{ik} \in \{0,1\} &\forall (i,k) \in N\times N \label{z-sbf}\\
&&& q_{k} \in \mathbb{R}^2 &\forall k \in N \label{q-sbf}\\
&&& t_{k} \geq 0 & \forall k \in N \label{t-sbf}
\end{align}
\end{subequations}

 The objective function of \ref{sbf} minimizes the total distance traveled, constraint \eqref{tour-sbf-1} ensures that each neighborhood is visited at some stage, while constraint \eqref{tour-sbf-2} makes it so that one neighborhood is visited at each stage. Constraint \eqref{anti-sym} fixes the first neighborhood to be visited, in order to eliminate symmetrical solutions given by shifting the sequence, constraint \eqref{dist-sbf} defines the distance between representative points (taking $n+1$ to be equal to $1$), and constraint \eqref{neigh-sbf} ensures that the representative point of stage $k$ is contained in the neighborhood visited in that stage. Finally, constraints \eqref{z-sbf}-\eqref{t-sbf} state the nature of the variables.



\section{MILP formulations}\label{sec4}

Given the inherent complexity associated with solving the previous problems due to their second-order cone constraints, and the superior performance of state-of-the-art solvers for MILP, it is useful to approximate these conic constraints with polyhedral representations. For this purpose, Ben-Tal and Nemirovski \cite{extended} propose a linear approximation of conic constraints such that the resulting feasible region is arbitrarily close to the original feasible region, with a limited number of variables and constraints.

Formally, let $\epsilon > 0$ and $\mathbf{L}^2 = \left\{(y,w) \in \mathbb{R}^2 \times \mathbb{R} \mid w \geq ||y||\right\}$ be the 3-dimensional Lorentz cone. We say that a polyhedron $P_\epsilon \subseteq \mathbb{R}^2 \times \mathbb{R}$ is an $\epsilon$-approximation of $\mathbf{L}^2$ if every point $(y,w) \in P_\epsilon$ is such that $||y||\leq (1+\epsilon)w$
.
In this context, Ben-Tal and Nemirovski \cite{extended} show that $P_\epsilon(\nu)$ is an $\epsilon(\nu)$-approximation of $\mathbf{L}^2$, where $\nu \in \mathbb{Z}_+$ is a construction parameter that determines the precision, $\epsilon(\nu)$ is given by

\begin{equation*}
    \epsilon(\nu) = \left(\cos\left(\frac{\pi}{2^{\nu + 1}}\right)\right)^{-1} - 1,
\end{equation*}

and the polyhedron $P_{\epsilon(\nu)}$ is defined by the following system of linear inequalities:

\begin{subequations}
\begin{align}
    P_\epsilon(\nu) := \bigg\{(&y,w) \in \mathbb{R}^2 \times \mathbb{R}:\\
    &\xi^0 \geq |y_1| \label{ext-first}\\
    &\eta^0 \geq |y_2| \\
    &\xi^j = \cos\left(\frac{\pi}{2^{j+1}}\right)\xi^{j-1} + \sin\left(\frac{\pi}{2^{j+1}}\right) \eta^{j-1} & \forall j \in \{1,\dots,\nu\}\\
    &\eta^j \geq \left| -\sin\left(\frac{\pi}{2^{j+1}}\right) \xi^{j-1} + \cos\left(\frac{\pi}{2^{j+1}}\right) \eta^{j-1}\right| & \forall j \in \{1,\dots,\nu\}\\
    &\xi^\nu \leq w \\
    &\eta^\nu \leq \tan\left(\frac{\pi}{2^{\nu + 1}}\right) \xi^\nu \label{ext-last}\\
    & \xi^j, \eta^j \in \mathbb{R} & \forall j \in \{0,\dots,\nu\} \bigg\} \label{ext-variables}
\end{align}
\end{subequations}  

Thus, the approximation is given by the projection onto the $(y,w)$-space of the polyhedron defined by the system of inequalities \eqref{ext-first}-\eqref{ext-variables}. We note that the number of constraints and new variables grow linearly with $\nu$, in particular, the system has $2(\nu + 1)$ additional variables and $2(\nu + 2)$ constraints. Moreover, small values of $\nu$ already yield close approximations (e.g. setting $\nu = 3$ yields $\epsilon(\nu) < 0.02$), thus the resulting polyhedron $P_{\epsilon(\nu)}$ remains manageable in terms of dimensionality.

Figure \ref{fig:approximation_cross_sections} ilustrates the evolution of the approximation as $\nu$ increases, showing cross-sections of the approximation $P_{\epsilon(\nu)}$ and of the Lorentz cone on $w = 1$, for different values of $\nu$. In each subfigure, the inner disk corresponds to the cross-section of the Lorentz cone, while the outer polygon corresponds to the approximation $P_{\epsilon(\nu)}$. This figure clearly shows that the approximation becomes closer to the Lorentz cone as $\nu$ increases, and that it is already quite close for small values of $\nu$.

\begin{figure}[h]
    \centering
    \includegraphics[width=\textwidth]{feasible_region_w1.png}
    \caption{Cross-sections of the Lorentz cone and its approximation for different values of $\nu$} \label{fig:approximation_cross_sections}
\end{figure}

We utilize this approximation to reframe both our previous SOCP models as MILPs, devoid of any conic constraints. In the arc case, the new MILP formulation (\ref{abf-approx}) is given by:

\begin{subequations}
\begin{align}
\min &&& \sum_{i \in N}\sum_{j\in N} d_{ij}^a & \tag{ABF-A} \label{abf-approx}\\
\text{s.t.} &&& \eqref{tour-abf-1}-\eqref{sec-4}, \eqref{big-m-abf}-\eqref{p-abf}& \nonumber  \\
&&& \left(p_i - c_i, r_i\right) \in P_{\epsilon(\nu)} & \forall i \in N \label{dist-labf}\\
&&& \left(p_i - p_j, d_{ij}\right) \in P_{\epsilon(\nu)} & \forall (i,j) \in N \times N \label{neigh-labf}
\end{align}
\end{subequations}
The objective function is the same as in the original formulation, the first constraint group states the unchanged constraints from the original model, and \eqref{dist-labf} along with \eqref{neigh-labf} ensure that the solution lies in the approximation of the starting conic constraints.

We apply this same approach to obtain the following new formulation (\ref{sbf-approx}) for the sequence case:
\begin{subequations}
\begin{align}
\min &&& \sum_{k \in N} t_k \tag{SBF-A} \label{sbf-approx}\\
\text{s.t.} &&& \eqref{tour-sbf-1}-\eqref{anti-sym}, \eqref{z-sbf}-\eqref{t-sbf}& \nonumber \\
&&& \left(q_k - q_{k+1}, t_k\right) \in P_{\epsilon(\nu)} & \forall k \in N \label{neigh-lsbf}\\
&&& \left(\sum_{i \in N} {z}_{ik}  c_i - q_k, \sum_{i\in N} {z}_{ik}  r_i\right) \in P_{\epsilon(\nu)} & \forall k \in N \label{dist-lsbf}
\end{align}
\end{subequations}
As in the previous formulation, the objective function remains unchanged from \ref{sbf}, the first constraint group states the non-conic constraints, and \eqref{neigh-lsbf} with \eqref{dist-lsbf} force the solution to be in the polyhedral approximation of the original feasible region induced by the conic constraints.

Since the approximation utilized is an outer-approximation, the feasible region obtained is slightly larger, and, as such, optimal solutions obtained could possibly not be feasible for the original problem, instead providing lower bounds. In order to obtain a feasible solution (and thus, an upper bound), one can take the visit sequence returned as optimal, fix the binary variables in the corresponding original formulation, and solve for the remaining variables with second-order cone programming.

\section{Decomposition schemes}\label{sec5}

In this section, we present decomposition schemes for all of our previously presented formulations. In the case of the SOCP formulations, these decompositions aim to reduce the complexity inherent to these kind of models, whereas in the case of the MILP approximations, the decompositions aim to close the optimality gap left by the approximation factor.

All four decompositions have the following structure: (i) a master problem that is a relaxation of the original problem and, as such, estimates the distance traveled; (ii) and a subproblem that computes the actual optimal distance for the visit sequence obtained in the master problem. In case the estimation obtained in the master problem differs from the value of the subproblem, a new cut is generated and added to the master problem. This procedure is iterated until the values obtained in the master problem and in the subproblem match for some solution.
\subsection{Decomposition of ABF}

First, we present the decomposition of \ref{abf} (\ref{abf-d}). In this decomposition, the master problem retains all tour-related variables and constraints, using underestimations of the distance between neighborhoods as weights in the objective function, while the subproblem takes the visit sequence given by the master problem, and computes the difference between the estimation and the actual optimal tour length for said sequence.

The master problem is as follows:
\begin{subequations} % placeholder
\begin{align}
\min &&& \sum_{i \in N}\sum_{j\in N} x_{ij} m_{ij} + \theta & \label{abf-d} \tag{ABF-D}\\
\text{s.t.} &&& \eqref{tour-abf-1}-\eqref{sec-4}, \eqref{x-abf}& \nonumber \\
&&& \theta \geq 0 & \label{theta-lb-abf-d}\\
&&& \theta \geq a^Tx + b & \forall (a,b) \in C \label{theta-estim-abf-d}
\end{align}
\end{subequations}
Here, $m_{ij}$ is an estimator of the distance between neighborhoods $i$ and $j$, for which we use the minimum possible distance between said neighborhoods. The new variable $\theta$ represents the difference between the estimation of the total distance and the actual distance for a given visit sequence. Set $C$ is originally empty, and contains all cuts added to the master problem. In this context, the first constraint group guarantees that the variable $x$ induces a valid tour, while \eqref{theta-lb-abf-d} sets a starting lower bound for $\theta$ (since we are underestimating the tour distance, this difference will always be non-negative, thus making 0 a valid lower bound), and \eqref{theta-estim-abf-d} ensures that the value of $\theta$ does not underestimate the difference between the estimated total distance and the actual distance. 

The subproblem is as follows:
\begin{subequations}
\begin{align}
Q^A(x) = \min &&& \sum_{(i,j) \in S(x)} \left(d_{ij} - m_{ij}\right) & \label{obj-subpr-abf}\\
\text{s.t.} &&& ||p_i - c_i|| \leq r_i &\forall i \in N \label{neigh-abf-d}\\
&&& ||p_i - p_j|| \leq d_{ij} &\forall (i,j) \in S(x) \label{dist-abf-d}\\
&&& d_{ij}\geq 0 &\forall (i,j) \in S(x) \label{d-abf-d}\\
&&& p_i \in \mathbb{R}^2 &\forall i \in N \label{p-abf-d}
\end{align}
\end{subequations}
In the subproblem, set $S(x) = \left\{(i,j) \in N \times N \mid x_{ij} = 1\right\}$ is the support of $x$, and, as such, the constraints ensure that the points selected for the visit sequence given by such $x$ satisfy neighborhood constraints and properly activate the distance variables, which are then minimized by the objective function. Note that, since we instantiate the subproblem over $S(x)$, which always has cardinality $n$, instead of over $N \times N$, the subproblem has $2n$ conic constraints, thus, a considerably lesser amount than the original $n^2 + n$ conic constraints in \ref{abf}.

 In order to solve this decomposition, we solve the master problem and, given an integer solution $\hat{x}$, we compute $Q^A(\hat{x})$. If the corresponding solution $\hat{\theta}$ is such that $\hat{\theta} < Q^A(\hat{x})$, then we add the following optimality cuts to the master problem \cite{ILSM}:
\begin{subequations}
\begin{gather}
     -\frac{1}{2}Q^A(\hat{x})\left(\sum_{(i,j) \in S(\hat{x})} (1 - x_{ij})\right) + Q^A(\hat{x}) \leq \theta \label{ll-cut}\\
     -\frac{1}{2}Q^A(\hat{x})\left(\sum_{(i,j) \in S(\hat{x})} (1 - x_{ji})\right) + Q^A(\hat{x}) \leq \theta \label{ll-cut-sim}
\end{gather}
\end{subequations}
Note that \eqref{ll-cut} exclusively eliminates the solution $\left(\hat{x},\hat{\theta}\right)$, since, for any feasible $x$ different from $\hat{x}$, they will differ in at least two entries, thus leading the cut to being dominated by constraint \eqref{theta-lb-abf-d}, and, as such, trivially satisfied. The cut \eqref{ll-cut-sim} does the same for the visit sequence given by reversing the arcs utilized, which will always induce the same total distance due to the symmetric nature of the problem.

Since these cuts are, as previously mentioned, enumerative in nature, it could potentially be necessary to enumerate every single possible tour and compute their corresponding distances in order to obtain the optimal solution.

\subsection{Decomposition of ABF-A}

Now we present the decomposition of \ref{abf-approx} (\ref{abf-a-d}). In this decomposition, the master problem simply consists of \ref{abf-approx} with an added variable that estimates the difference between the distance obtained by the approximation and the actual distance, while the subproblem computes the value of said difference for a given visit sequence.

The master problem is as follows:
\begin{subequations}
\begin{align}
\min &&& \sum_{i \in N}\sum_{j\in N} d_{ij}^a +\theta & \tag{ABF-A-D} \label{abf-a-d}\\
\text{s.t.} &&& \eqref{tour-abf-1}-\eqref{sec-4}, \eqref{big-m-abf}-\eqref{p-abf}& \nonumber \\
&&& \eqref{dist-labf}-\eqref{neigh-labf}  & \nonumber\\
&&& \theta \geq 0 & \label{theta-lb-abf-ext-d}\\
&&& \theta \geq a^Tx + b & \forall (a,b) \in C & \label{theta-estim-abf-ext-d}
\end{align}
\end{subequations}
As in the previous decomposition (\ref{abf-d}), variable $\theta$ represents the difference between the estimated distance and the real optimal distance, while $C$ is the set of all optimality cuts added. The constraints are the same as in \ref{abf-approx}, adding an initial lower bound for the estimator $\theta$ in \eqref{theta-lb-abf-ext-d}, and the set of optimality cuts in \eqref{theta-estim-abf-ext-d}.
The subproblem is as follows:
\begin{equation}
Q^A_\nu(x, d^a) = \min \left\{ \sum_{(i,j) \in S(x)} \left(d_{ij} - d_{ij}^a\right) : \eqref{neigh-abf-d}-\eqref{p-abf-d} \right\} 
\end{equation}
This subproblem is the same as the \ref{abf-d} subproblem, just changing the objective function to adapt to the estimation $d^a$ returned by the master problem. Just as before, we solve the master problem and, once a solution $(\hat{x}, \hat{d^a}, \hat{\theta})$ is found, we solve $Q^A_\nu(\hat{x},\hat{d^a})$ and get its optimal value. If $\hat{\theta} < Q^A_\nu(\hat{x},\hat{d^a})$, we add cuts \eqref{ll-cut} and \eqref{ll-cut-sim}, replacing $Q^A(\hat{x})$ by $Q^A_\nu(\hat{x},\hat{d^a})$.

\subsection{Decomposition of SBF}

Now we present the decomposition of \ref{sbf} (\ref{sbf-d}). In this decomposition, the master problem determines which neighborhood gets visited at each stage, while the subproblem computes the distance associated with said visit order. The master problem is as follows:
\begin{subequations} 
\begin{align}
\min &&& \theta & \tag{SBF-D} \label{sbf-d}\\
\text{s.t.} &&& \eqref{tour-sbf-1}-\eqref{anti-sym}, \eqref{z-sbf}& \nonumber \\
&&& \theta \geq 0 & \label{theta-lb-sbf-d}\\
&&& \theta \geq a^Tz + b & \forall (a,b) \in C & \label{theta-estim-sbf-d}
\end{align}
\end{subequations}
Here, the variable $\theta$ estimates the total distance of the tour determined by variable $z$, with the first group of constraints ensuring that said tour is valid, while \eqref{theta-lb-sbf-d} and \eqref{theta-estim-sbf-d} work just as in the previous decompositions. 
The subproblem for this decomposition is as follows:
\begin{subequations}
\begin{align}
Q^S(z) = \min &&& \sum_{k \in N} t_k  & \label{obj-subpr-sbf}\\
\text{s.t.} &&& ||q_k - q_{k+1}|| \leq t_k &\forall k \in N \label{dist-sbf-d}\\
&&& \left|\left|\sum_{i \in N} {z}_{ik}  c_i - q_k\right|\right| \leq \sum_{i\in N} {z}_{ik}  r_i &\forall k \in N \label{neigh-sbf-d}\\
&&& q_{k} \in \mathbb{R}^2 &\forall k \in N \label{q-sbf-d}\\
&&& t_{k} \geq 0 & \forall k \in N \label{t-sbf-d}
\end{align}
\end{subequations}
The constraints in this subproblem are identical to those that determine the visit points and the distances in \ref{sbf}, thus, it computes the optimal distance for the visit sequence encoded by $z$. Since the objective function of the master problem is simply $\theta$, and the cuts we use for \ref{abf-d} are of an enumerative nature, using these same cuts here would force us to effectively iterate over every single possible sequence, rapidly rendering them ineffective as $n$ grows. 

In order to deal with this issue, we can take the dual of $Q^S(z)$, and use duality theory to formulate cuts as needed. Note that the subproblem satisfies Slater's condition and is bounded from below 
(by 0), so strong duality holds \cite{slaters-cond}. Taking the dual, we can re-formulate $Q^S(z)$ as follows:
\begin{subequations} \label{dual-sbf}
\begin{align}
Q^S(z) = \max &&& \sum_{i \in N}\sum_{k \in N} \mu_{ik} z_{ik} & \label{obj-dual-sbf}\\
\text{s.t.} &&& 1 - \lambda_k^n - \lambda_k^d = 0 & \forall k \in N\\
&&& \mu_{ik} + c_i \cdot \rho^c_k + r_i\lambda^c_k = 0 &\forall (i,k) \in N \times N \label{mu-con-sbf-dual}\\
&&& \rho_k^d - \rho_{k-1}^d + \rho_k^c = 0 & \forall k \in N \label{order-con}\\
&&&||\rho_k^d|| \leq \lambda_k^d & \forall k \in N\\
&&& ||\rho_k^c|| \leq \lambda_k^c & \forall k \in N\\
&&& \lambda_k^n, \lambda_k^d, \lambda_k^c \geq 0 & \forall k \in N \\
&&& \mu_{ik} \in \mathbb{R} & \forall (i,k) \in N \times N\\
&&& \rho_k^d, \rho_k^c \in \mathbb{R}^2 & \forall k \in N
\end{align}
\end{subequations}
Variables with superscript $d$ are dual variables associated with distance constraints \eqref{dist-sbf-d}, variables with superscript $c$ are those associated with distance to the center constraints \eqref{neigh-sbf-d}, and variables with superscript $n$ are associated with variable nature constraints \eqref{t-sbf-d}. For brevity in constraint \eqref{order-con}, we assume $k - 1 = n$ when $k = 1$. Finally, $\mu$ is associated with an auxiliary constraint utilized to build the dual problem, and can in fact be eliminated, since it is fully determined by constraint \eqref{mu-con-sbf-dual}, so we can integrate this constraint into the objective function, eliminating the variable $\mu$. However, for notational simplicity, we retain this variable in the model.

Now, given an integer solution $\hat{z}$ to the master problem, we compute $Q^S(\hat{z})$, and, if $\hat{\theta} < Q^S(\hat{z})$, we retrieve the solution $\hat{\mu}$ of the subproblem and add the following optimality cut:
\begin{equation}
    \sum_{i \in N}\sum_{k \in N} \hat{\mu}_{ik} z_{ik} \leq \theta
\end{equation}

\begin{remark}
    We could use this dual approach in \ref{abf-d}, however, we opt not to because variable $x$ does not appear as a parameter in the subproblem $Q^A(x)$. Thus, in order to be able to obtain similar cuts, we would need to reformulate the subproblem by instantiating over $N \times N$ instead of $S(x)$ and carrying over the big-M constraint \eqref{big-m-abf} to make $x$ explicit in the model, thus considerably increasing the dimension of the model. Additionally, the cut obtained involves the big-M parameter and, in our preliminary experiments, lead to worse results than simply utilizing the proposed cuts \eqref{ll-cut} and \eqref{ll-cut-sim}.
\end{remark}

\subsection{Decomposition of SBF-A}

Following the same principle as in \ref{abf-a-d}, we present the decomposition of \ref{sbf-approx} (\ref{sbf-a-d}), whose master problem simply adds a difference estimator $\theta$, along with the corresponding constraints, to \ref{sbf-approx}, resulting in:
\begin{subequations} 
\begin{align}
\min &&& \sum_{k \in N} t_k +\theta & \tag{SBF-A-D} \label{sbf-a-d}\\
\text{s.t.} &&&\eqref{tour-sbf-1}-\eqref{anti-sym}, \eqref{z-sbf}-\eqref{t-sbf}& \nonumber  \\
&&& \eqref{neigh-lsbf}-\eqref{dist-lsbf} & \nonumber \\
&&& \theta \geq 0 & \label{theta-lb-sbf-ext-d}\\
&&& \theta \geq a^Tz + b & \forall (a,b) \in C & \label{theta-estim-sbf-ext-d}
\end{align}
\end{subequations}
As a subproblem for this decomposition, we define $Q_\nu^S(\hat{z},\hat{t}) = Q^S(\hat{z}) - \sum_{k\in N} \hat{t}_k$. We can use either the primal or the dual version of the subproblem utilized in \ref{sbf-d}, $Q^S(\hat{z})$. In both cases, if we get a solution such that $\hat{\theta} < Q^S_\nu(\hat{z})$, we add the following cuts, analogous to the ones used in \ref{abf-d}:
\begin{subequations}
\begin{gather}
        -\frac{1}{2}Q^S_\nu(\hat{z})\left(\sum_{i,k \in S(\hat{z})} (1 - z_{ik})\right) + Q^S_\nu(\hat{z}) \leq \theta \label{ll-cut-sbf}\\
     -\frac{1}{2}Q^S_\nu(\hat{z})\left(\sum_{i,k \in S(\hat{z})} (1 - z_{i(n-k+1)})\right) + Q^S_\nu(\hat{z}) \leq \theta \label{ll-cut-sim-sbf}
\end{gather}
\end{subequations}
Here, \eqref{ll-cut-sbf} works just as \eqref{ll-cut}, and \eqref{ll-cut-sim-sbf} is the cut corresponding to the sequence obtained by reversing the visit order described by $\hat{z}$, which has the same optimal tour distance, enabling us to include this cut. Alternatively, if using the dual version of the subproblem, we can add the following cut:
\begin{equation}
    \sum_{i \in N}\sum_{k \in N} \hat{\mu}_{ik} z_{ik} - \sum_{k \in N} \hat{t}_k \leq \theta \label{dual-cut-sbf}
\end{equation}
This cut is analogous to the one used in \ref{sbf-d}, just adapting it to account for the additional constant term in the subproblem's objective function.


\section{On the applicability of generalized Benders decomposition}\label{sec6}


The seminal work of \citet{Benders} introduces a decomposition procedure for the solution of MILP problems with \textit{complicating variables}, i.e., variables that, when fixed, yield a simpler (LP) subproblem, for which linear duality theory can be applied to develop cuts. This procedure, known as Benders Decomposition, has been widely used in the literature for the solution of large-scale MILP problems \cite{benders-survey}. 

While highly effective for MILP problems, Benders decomposition (at least in its original form) is limited to problems where the subproblem takes the form of a continuous linear formulation. To extend the applicability of this decomposition technique to problems with non-linear subproblems, \citet{gbd} proposes Generalized Benders Decomposition (GBD), which allows for the decomposition of problems where the subproblem is a convex optimization problem, with applications on some non-convex problems as well \cite{gbd-applications}.

GBD considers problems of the following form:
\begin{equation}
    \min_{x,y} \left\{f(x,y) \text{ s.t. } G(x,y) \leq 0, x\in X, y \in Y\right\}
\end{equation}
where $X$ and $Y$ are some real-valued domains, $f: X \times Y \rightarrow \mathbb{R}$ is the objective function, and $G: X \times Y \rightarrow \mathbb{R}^m$ is a vector of constraint functions. This problem can be re-stated as
\begin{equation}
    \min_{x \in X} \left\{v(x) \text{ s.t. } x \in X\cap V\right\},
\end{equation}
where
\begin{equation}
    v(x) = \inf_{y\in Y} \left\{f(x,y) \text{ s.t. } G(x,y) \leq 0\right\}
\end{equation}
and
\begin{equation}
    V = \left\{x \in X: G(x,y) \leq 0 \text{ for some } y\in Y\right\}.
\end{equation}
Under some -- fairly weak -- assumptions on $f$, $G$, and $Y$ \cite{gbd}, it holds that
\begin{equation}
    v(x) = \sup_{u \geq 0} \left[\inf_{y\in Y} \left\{f(x,y) + u\transpose G(x,y)\right\}\right].
\end{equation}
With this result, the original problem can be reformulated as
\begin{align}
    \min_{x\in X}\quad & \theta \\
    \text{s.t.} \quad& \theta \geq \inf_{y\in Y} \left\{f(x,y) + u\transpose G(x,y)\right\}, \quad \forall u \geq 0  \label{gbd-opt}\\
    & x \in V \label{gbd-feas}
\end{align}
The logic of the decomposition is to relax constraints \eqref{gbd-opt} and \eqref{gbd-feas} to obtain a master problem in variable $x$, and, given a solution $\hat{x}$ to the master problem, solve the subproblem $v(\hat{x})$, and add optimality cuts (from \eqref{gbd-opt}) or feasibility cuts (from \eqref{gbd-feas}) as needed, iterating until convergence. In the case of the CETSP, the subproblem is always feasible, so we focus solely on optimality cuts. These cuts take the form
\begin{equation}\label{gbd-cut}
    L^*(x, \hat{u}) \leq \theta,
\end{equation}
where
\begin{equation}
    L^*(x, \hat{u}) = \inf_{y\in Y} \left\{f(x,y) + \hat{u}\transpose G(x,y)\right\},
\end{equation}
and $\hat{u}$ is the optimal dual vector associated with the subproblem $v(\hat{x})$.

An important remark regarding the viability of GBD is that, in order for the cuts \eqref{gbd-cut} to be practical, it is necessary that the function $L^*(\cdot, \hat{u})$ on $X$ can be obtained explicitly in a way that requires little or no more effort than that which is required to evaluate it at a single value of $x$, which the author refers to as Property (P). Otherwise, the functional form of the cuts would be too complex to be of any practical use. An example where Property (P) holds is when both $f$ and $G$ are linearly separable in $x$ and $y$. 

\citet{bagajewicz} show that, if we assume that for each $\hat{u} \geq 0$, there exists some $\hat{y} \in Y$ such that the cuts can be expressed as
\begin{equation}
    L^*(x,\hat{u}) = \inf_{y\in Y} \left\{f(x,y) + \hat{u}\transpose G(x,y)\right\} = f(x,\hat{y}) + \hat{u}\transpose G(x,\hat{y}),
\end{equation}
which, as a condition, is stronger than Property (P), then it is possible for the GBD procedure to converge to a non-optimal solution that can, in fact, not even be a local minimum.

% if a slightly stronger version of Property (P) is assumed in order to obtain an explicit form for the cuts, it is possible for the procedure to converge to a point that is not even a local minimum of the original problem, thus highlighting the importance of this property and proper formulation of the cuts.

% =========================================================================================

In the context of the PDTSPN, \citet{walteros} propose a formulation similar to \ref{abf}, for which they implement a decomposition scheme based on GBD. The resulting master problem in their decomposition is the same as \ref{abf-d} (with additional constraints corresponding to the pickup-and-delivery aspect of the problem), and the subproblem is as follows:
\begin{subequations}
\begin{align}
Q(x) = \min &&& \sum_{i,j \in N, i \neq j} x_{ij} (d_{ij} - m_{ij})\\
\text{s.t.} &&& f_i(p_i) \leq 0 &\forall i \in N \label{neigh-constraint-pdtspn}\\
&&& ||p_i - p_j|| \leq d_{ij} &\forall (i,j) \in N \times N \label{dist-constraint-pdtspn}\\
&&& d_{ij}\geq 0 &\forall (i,j) \in N \times N \\
&&& p_i \in \mathbb{R}^2 &\forall i \in N
\end{align}
\end{subequations}
where functions $f_i$ represent the neighborhoods, which are assumed to be convex and compact. In particular, circular neighborhoods can be represented as a specific case of these functions.

This subproblem differs from the one we present in \ref{abf-d} in that it is instantiated over all arcs in $N \times N$, instead of only over the arcs in the support of $x$, $S(x)$. 

For simplicity, we assume that the distance estimations $m_{ij}$ are all equal to 0. In this context, cuts \eqref{gbd-cut} take the following form:
\begin{equation}
    \underbrace{\inf_{(p, d) \in P }\left\{\sum_{i,j \in N, i \neq j} d_{ij}x_{ij} - \sum_{i \in N} \hat{\pi}_i f_i(p_i) - \sum_{(i,j) \in N \times N}\hat{\mu}_{ij} (||p_i - p_j|| - d_{ij})\right\}}_{L^*(x,\hat{\pi}, \hat{\mu})} \leq \theta
\end{equation}
where 
\begin{align*}
    P= \{(p,d) : & f_i(p_i) \leq 0 \quad &\forall i \in N, \\
    & ||p_i - p_j|| \leq d_{ij} \quad &\forall (i,j) \in N \times N, \\
    & d_{ij} \geq 0 \quad &\forall (i,j) \in N \times N, \\
    & p_i \in \mathbb{R}^2 \quad &\forall i \in N \}
\end{align*}
is the feasible region of the subproblem (note that it does not depend on $x$), and $\hat{\pi}$ and $\hat{\mu}$ are the optimal dual vectors associated with the neighborhood constraints \eqref{neigh-constraint-pdtspn} and the distance constraints \eqref{dist-constraint-pdtspn}, respectively.

Based on this formulation, \citet{walteros} propose that, given a master problem solution $\hat{x}$, and the corresponding primal and dual optimal solutions of the subproblem $(\hat{p}, \hat{d})$ and $(\hat{\pi}, \hat{\mu})$, respectively, the following cut can be added to the master problem if needed:
\begin{equation}
    \sum_{i,j \in N, i \neq j} \hat{d}_{ij} x_{ij} - \sum_{i \in N} \hat{\pi}_i f_i(\hat{p}_i) - \sum_{(i,j) \in N \times N}\hat{\mu}_{ij} (||\hat{p}_i - \hat{p}_j|| - \hat{d}_{ij}) \leq \theta
\end{equation}
Since $x$ does not appear in the second two summations, by complementary slackness, these terms are equal to zero, thus simplifying the cut to
\begin{equation}
    \sum_{i,j \in N, i \neq j} \hat{d}_{ij} x_{ij} \leq \theta
\end{equation}
However, this cut is not valid, since
\begin{equation}
    L^*(x,\hat{\pi}, \hat{\mu}) \leq \sum_{i,j \in N, i \neq j} \hat{d}_{ij} x_{ij} - \sum_{i \in N} \hat{\pi}_i f_i(\hat{p}_i) - \sum_{(i,j) \in N \times N}\hat{\mu}_{ij} (||\hat{p}_i - \hat{p}_j|| - \hat{d}_{ij}),
\end{equation}
as $(\hat{p}, \hat{d})$ is a feasible solution to the problem defining $L^*(x,\hat{\pi}, \hat{\mu})$, but not necessarily optimal. Thus, the cuts proposed are overestimating $L^*(x,\hat{\pi}, \hat{\mu})$, and, as such, are not valid. From the results of \citet{bagajewicz}, it is possible that the use of these cuts could lead to convergence to a non-optimal solution.

We will now present two proofs for the invalidity of these cuts, the first one related to numerical issues that arise on implementation, and the second one being a concrete example. First, let us consider the sequence $\tilde{S}$ associated with some first-stage solution $\tilde{x}$, i.e. $\tilde{S} = \left\{(i,j) \in N \times N \mid \tilde{x}_{ij} = 1\right\}$. Note that, from $x$ being a binary variable, the objective function of the subproblem can be rewritten as
\begin{equation}
    \sum_{i,j \in N, i \neq j} d_{ij} \tilde{x}_{ij} = \sum_{(i,j) \in \tilde{S}} d_{ij}
\end{equation}
Let us note by $\tilde{d}$ the corresponding vector of optimal distances obtained by solving the subproblem. Note that, since only entries $i,j \in \tilde{S}$ appear in the objective function, and there is no upper bound on any of the distance variables, the optimal distance $\tilde{d}_{ij}$ for $i,j \notin \tilde{S}$ can take any non-negative value greater or equal to $||p_i - p_j||$, without affecting the optimality of the solution. Thus, the entries $\tilde{d}_{ij}$ for $i,j \notin \tilde{S}$ are not necessarily defined as the distance between the corresponding visit points, and can in fact be arbitrarily large, potentially leading to cuts that incorrectly eliminate feasible solutions.

We now present a numerical example that illustrates how these interpretation of \eqref{gbd-cut} can cut off feasible solutions. Consider an instance with four nodes, whose coordinates and radii are given in Table \ref{tab:counterexample_inst}, and plotted in Figure \ref{fig:counterexample_inst}.


\begin{figure}[htbp]
\centering

% Left: figure content
\begin{minipage}{0.49\linewidth}
    \centering
    \includegraphics[width=\linewidth]{counterexample_instance_plot.png}

    \captionof{figure}{Plot of the example instance}
    \label{fig:counterexample_inst}
\end{minipage}
\hfill
% Right: two tables stacked
\begin{minipage}{0.49\linewidth}
    \centering

    %--- Table 1 ---
    \begin{tabular}{cccc}
    \toprule
    Node & x & y & r \\
    \midrule    
    0 & 3 & 2 & 0 \\
    1 & 5 & 5 & 0.5\\
    2 & 4 & 2 & 0.25 \\
    3 & 2 & 8 & 0.25 \\
    \bottomrule
    \end{tabular}
    \captionof{table}{Coordinates and radii}
    \label{tab:counterexample_inst}

    \vspace{1.5em}

    %--- Table 2 ---
    \begin{tabular}{c|cccc}
         & 0 & 1 & 2 & 3 \\ \hline
    0    & 0.00 & 3.11 & 1.03 & 5.83 \\
    1    & 3.11 & 0.00 & 2.43 & 4.20 \\
    2    & 1.03 & 2.43 & 0.00 & 5.84 \\
    3    & 5.83 & 4.20 & 5.84 & 0.00 \\
    \end{tabular}
    \captionof{table}{Distance matrix for $0 \rightarrow 1 \rightarrow 2 \rightarrow 3 \rightarrow 0$}
    \label{tab:counterexample_dist_0123}

\end{minipage}

\end{figure}





The optimal solution to this instance has a total distance of approximately 13.24, and the optimal visit order is $0 \rightarrow 3 \rightarrow 1 \rightarrow 2 \rightarrow 0$ (or the reverse). If we solve the subproblem for the visit order $0 \rightarrow 1 \rightarrow 2 \rightarrow 3 \rightarrow 0$, we can compute the distances manually given the optimal visit points to avoid the problem mentioned before, thus obtaining the distance matrix presented in Table \ref{tab:counterexample_dist_0123}.


Let us denote the optimal sequence by $S^*$, the second sequence by $\tilde{S}$, and its corresponding previously presented distance matrix by $\tilde{d}$. Suppose we generated the cut corresponding to $\tilde{S}$:
\begin{equation}
    \sum_{i,j \in N, i \neq j} x_{ij} \tilde{d}_{ij} \leq \theta
\end{equation}

which, given the distance matrix in Table \ref{tab:counterexample_dist_0123}, can be rewritten approximately as
\begin{equation}
\begin{aligned}
&3.11(x_{01} + x_{10}) + 1.03(x_{02} + x_{20}) + 5.83(x_{03} + x_{30}) \\
&+ 2.43(x_{12} + x_{21}) + 4.20(x_{13} + x_{31}) + 5.84(x_{23} + x_{32})
\leq \theta
\end{aligned}
\end{equation}

Now, if we evaluate the left-hand side of this cut for sequence $S^*$, we obtain
\begin{align*}
    \sum_{i,j \in S^*} \tilde{d}_{ij} &= \tilde{d}_{03} + \tilde{d}_{31} + \tilde{d}_{12} + \tilde{d}_{20}\\
    &\approx 5.83 + 4.20 + 2.43 + 1.03 = 13.49 
\end{align*}
which is greater than the optimal distance of approximately 13.24.

Thus, the cut obtained by solving the subproblem for sequence $\tilde{S}$ leads to an overestimation of the distance of the optimal sequence $S^*$, and is therefore incorrect. This example illustrates that the cuts proposed by \citet{walteros} are invalid, and, as such, the decomposition scheme they propose does not guarantee convergence to the optimal solution of the CETSP, or its extension, the PDTSPN.

\begin{remark}
It is important to note that we have not been able to verify whether Property (P) holds for the subproblem associated to the CETSP, and, as such, we cannot categorically state that GBD is not applicable to this problem. However, the invalidity of the cuts proposed by \citet{walteros} highlights the challenges associated with applying GBD to this context, and suggests that further research is needed to explore the applicability of GBD to the CETSP and related problems.
\end{remark}


\section{Computational experiments}\label{sec7}

In this section, we conduct computational experiments to test the different formulations presented in this work. We generate several CETSP instances via the following procedure. Given parameters $n$ for the number of points, $r$ for the mean radii, and $\sigma$ for the variation among radii, the coordinates of the $n$ nodes are chosen randomly from a uniform distribution on a square with side 10. Then, for each node, a radius is sampled uniformly from $[(1-\sigma)r, (1+\sigma)r]$. For each parameter combination, we generate five different instances. We perform this procedure over every combination of parameters $n \in \{10, 15, 20\}$, $r \in \{0.25, 0.5, 1\}$, and $\sigma \in \{0, 0.2, 0.5\}$, obtaining a total of 135 distinct instances. Approximations were implemented with $\nu = 3$, which yields an approximation factor of $\epsilon(\nu) < 0.02$.

Before solving, we consider a pre-processing step, where, if there is a pair $i,j \in N$ such that the neighborhood of $i$ is wholly contained in the neighborhood of $j$, we eliminate node $j$ from the problem, since it is trivially visited by any tour that visits the neighborhood of $i$.

We implement all formulations in Python 3.10, using Gurobi 12.0.2 as the solver, with a ten-minute time limit. All experiments are conducted on a PC equipped with an Intel Core i5-11300H 3.10GHz processor and 16GB of memory, running Windows 11. In the case of \ref{sbf-a-d}, we report results using cuts \eqref{ll-cut-sbf} and \eqref{ll-cut-sim-sbf}, as they yielded slightly better results compared to using \eqref{dual-cut-sbf}. A complete comparison of the two cut types is presented in appendix \ref{appendix}.

\subsection{General results for exact formulations} \label{gen-res}

We now review the results of the six different exact formulations over the 135 instances of the problem. Figure \ref{perf-profile-gen} presents a comparison of the performance of the six formulations as a function of execution time and the optimality gap. On the X-axis, values to the left of the vertical line represent time, while values to the right correspond to the optimality gap. The Y-axis indicates the number of instances solved. Each colored line represents the performance of a different formulation. The first vertical dotted line marks the 600 seconds time limit; thus, the left part of the plot shows the number of instances solved to optimality and the time needed to do so, while the right part of the plot presents information about the gap left on instances not solved to optimality.

Arc-based formulations are plotted in red, while sequence-based formulations are in blue. In both cases, the lighter shade corresponds to the base formulation, the medium shade to the decomposition, and the darker shade to the approximation with decomposition. 

\begin{figure}[H]
\centering
\includegraphics[width=0.6\linewidth]{performance_plot.png}
\caption{Performance Profile Plot of exact formulations on all instances}\label{perf-profile-gen}
\end{figure}

We observe that the formulations that solve the greater amount of instances to optimality are the two-stage formulations based on the approximation of the second-order cone, \ref{sbf-a-d} and \ref{abf-a-d}, with the first one solving slightly more instances. These are followed by \ref{abf-d}, and there is a big gap separating these three formulations from the other ones. We can also observe that, besides \ref{sbf-a-d}, all formulations find most of their optimal solutions in a fairly quick manner, and after about 300 seconds the amount of optimal solutions stagnates.

Table \ref{tab:general performance} supports these results, where the ``Formulation'' column states the formulation used, the ``\#Opt'' column indicates the number of optimal solutions found within the time limit, the ``TimeOpt'' column reports the average time on instances solved to optimality, while ``Gap'' and ``Gap (NoOpt)'' present the average gap across all instances, and the average gap without considering instances solved to optimality, respectively. Finally, the ``\#Cuts'' column shows the average number of cuts generated.

\begin{table}[h]
\centering
\begin{tabular}{lrrrrr}
\toprule
Formulation & \#Opt & TimeOpt & Gap & Gap (NoOpt) & \#Cuts \\
\midrule
ABF & 29 & 17.42 & 0.43 & 0.54 & - \\
ABF-D & 66 & 21.87 & 0.16 & 0.30 & 12311.35 \\
ABF-A-D & 84 & 74.18 & 0.12 & 0.33 & 109.27 \\
SBF & 43 & 46.29 & 0.43 & 0.63 & - \\
SBF-D & 45 & 73.06 & 0.60 & 0.89 & 2232.52 \\
SBF-A-D & 86 & 119.36 & 0.18 & 0.51 & 599.11 \\
\bottomrule
\end{tabular}
\caption{Summary of results for all instances}
\label{tab:general performance}
\end{table}

We observe that, overall, \ref{abf-a-d} has the best performance, with only two less instances solved to optimality than \ref{sbf-a-d}, but with a considerably smaller average gap, both overall and on instances not solved to optimality. With regards to the second and third formulations with best performance, \ref{sbf-a-d} and \ref{abf-d}, we note that \ref{sbf-a-d} solves 20 more instances to optimality, but with substantially larger average gaps, and \ref{abf-d} has smaller average runtime on optimal instances. Additionally, we note that, while \ref{abf} and \ref{sbf} cannot be categorically ranked, with \ref{sbf} achieving a greater number of optimalities and \ref{abf} smaller gaps on non-optimal instances, \ref{abf-d} seems to be a considerably better formulation than its sequence-based counterpart, obtaining smaller gaps and more optimalities in less average time. In fact, \ref{sbf-d} appears to be a fairly weak decomposition, obtaining only two more optimalities than \ref{sbf}, but with quite larger gaps on instances not solved to optimality.

It is also interesting to note that both non approximation-based arc formulations have average running times on optimal instances under 22 seconds, suggesting that, at least with the time limit imposed, these formulations either find an optimal solution quickly, or fail to find an optimal solution altogether. 

The superior performance of the decompositions based on the approximation of the second-order cone, \ref{sbf-a-d} and \ref{abf-a-d}, can be explained by the higher quality of the distance estimation used in the master problem. These formulations leverage the quality of the approximation of the conic constraints to obtain fairly accurate cost estimations, thus lessening the need for refinement via cut generation. On the other hand, the cost estimations used in \ref{abf-d} are fairly lax, and have the potential to severely underestimate the actual costs, leading to the need for a large number of cuts to be generated in order to obtain good solutions due to the enumerative nature of the cuts used. This is also a possible reason for the poor performance of \ref{sbf-d}, since it has no proper estimation of costs in the master problem, thus obtaining all information about the length of the tour from the cuts generated in the subproblem.

This is supported by the number of cuts generated by each formulation, where \ref{abf-a-d} and \ref{sbf-a-d} generate two orders of magnitude fewer cuts than \ref{abf-d}. This suggests that the quality of the master problem cost estimation is a key factor in the performance of these formulations, as it can significantly reduce the need for cut generation, which can be computationally expensive and time-consuming.

Lastly, it is worth remarking that the three best performing formulations all utilize cuts that are enumerative in nature, which suggests that there is still room for improvement in this regard, as stronger cuts could lead to considerably better results.

\subsection{Results by instance type for exact formulations} \label{spec-res}

To provide a more in-depth analysis of performance across different instance types, we present similar plots and tables as in \ref{gen-res}, while grouping instances according to the number of points, and to the value of the parameter $r$ used in the generation of the instances. We do not report the analysis grouping by value of $\sigma$, since taking different values of this parameter did not lead to significant performance variations, at least with the values tested.

\begin{figure}[H]
\centering
\includegraphics[width=0.6\linewidth]{performance_n_values.png}
\caption{Performance Profile Plot of exact formulations grouped by amount of points}\label{fig:profile-nodes}
\end{figure}

\begin{table}[h]
\centering
\begin{tabular}{l l r r r r r}
\toprule
n & Formulation & \#Opt & TimeOpt & Gap & Gap (NoOpt) & \#Cuts \\
\midrule
10 & ABF & 28 & 17.88 & 0.05 & 0.13 & - \\
10 & ABF-D & 45 & 11.00 & 0.00 & - & 1433.38 \\
10 & ABF-A-D & 45 & 5.44 & 0.00 & - & 51.78 \\
10 & SBF & 43 & 46.29 & 0.03 & 0.59 & - \\
10 & SBF-D & 45 & 73.06 & 0.00 & - & 1656.38 \\
10 & SBF-A-D & 45 & 5.79 & 0.00 & - & 201.96 \\
15 & ABF & 1 & 4.35 & 0.49 & 0.50 & - \\
15 & ABF-D & 20 & 28.89 & 0.17 & 0.30 & 15485.91 \\
15 & ABF-A-D & 27 & 92.07 & 0.11 & 0.26 & 136.18 \\
15 & SBF & 0 & 0.00 & 0.45 & 0.45 & - \\
15 & SBF-D & 0 & 0.00 & 0.79 & 0.79 & 3030.91 \\
15 & SBF-A-D & 41 & 244.02 & 0.01 & 0.15 & 1301.64 \\
20 & ABF & 0 & 0.00 & 0.74 & 0.74 & - \\
20 & ABF-D & 1 & 370.72 & 0.30 & 0.31 & 20014.76 \\
20 & ABF-A-D & 12 & 291.72 & 0.26 & 0.36 & 139.87 \\
20 & SBF & 0 & 0.00 & 0.81 & 0.81 & - \\
20 & SBF-D & 0 & 0.00 & 1.00 & 1.00 & 2010.27 \\
20 & SBF-A-D & 0 & 0.00 & 0.54 & 0.54 & 293.73 \\
\bottomrule
\end{tabular}
\caption{Summary of results for instances grouped by amount of points}
\label{tab:general performance nodes}
\end{table}

Figure \ref{fig:profile-nodes} and Table \ref{tab:general performance nodes} present the results grouped by the number of nodes in the instances. 

On instances with 10 points, \ref{abf-a-d} has the best performance, followed closely by \ref{sbf-a-d} and \ref{abf-d}, with all three of these formulations solving the whole set of instances to optimality in under 12 seconds. We also observe that \ref{sbf} has a better overall performance than \ref{abf} on these instances, only failing to achieve optimality on two of them, although \ref{abf} finds optimal solutions faster when it does find them. Lastly, it is interesting to note that \ref{sbf-d} solves all instances of this set to optimality. 

On instances with 15 points, the most notable feature is the performance of \ref{sbf-a-d}, which solves all but 4 of the 45 instances to optimality, and achieves very small gaps on those not solved, a substantially better performance than all other formulations. Once again, the second best performer is \ref{abf-a-d} followed closely by \ref{abf-d}. In contrast to the case with 10 points, here \ref{sbf-d} fails to solve any instance to optimality, and, in fact, is categorically outperformed by all other formulations. Finally, \ref{abf} has a similar performance to its sequence-based counterpart, solving one instance to optimality, but having larger gaps on average.

Finally, on instances with 20 points, we observe a considerable drop in performance across all formulations, with only \ref{abf-d} and \ref{abf-a-d} being able to solve any instance to optimality, with the latter being the best performer, solving 12 of the 45 instances. We also note that \ref{sbf-a-d}, which had been the best performer on instances with 15 points, fails to solve any instance to optimality here, and has a considerably larger average gap than \ref{abf-a-d}. Finally, we observe that \ref{abf} and \ref{sbf} have similar performances, with \ref{abf} having a slightly smaller average gap, and \ref{sbf-d} not only fails to solve any instance to optimality, but also has a gap of 100\% on all instances.

\begin{figure}[H]
\centering
\includegraphics[width=0.6\linewidth]{performance_r_values.png}
\caption{Performance Profile Plot of exact formulations grouped by radius values}\label{fig:profile-r-values}
\end{figure}

\begin{table}[h]
\centering
\begin{tabular}{l l r r r r r}
\toprule
r & Formulation & \#Opt & TimeOpt & Gap & Gap (NoOpt) & \#Cuts \\
\midrule
0.25 & ABF & 28 & 17.88 & 0.05 & 0.13 & - \\
0.25 & ABF-D & 45 & 11.00 & 0.00 & - & 1433.38 \\
0.25 & ABF-A-D & 45 & 5.44 & 0.00 & - & 51.78 \\
0.25 & SBF & 43 & 46.29 & 0.03 & 0.59 & - \\
0.25 & SBF-D & 45 & 73.06 & 0.00 & - & 1656.38 \\
0.25 & SBF-A-D & 45 & 5.79 & 0.00 & - & 201.96 \\
0.50 & ABF & 1 & 4.35 & 0.49 & 0.50 & - \\
0.50 & ABF-D & 20 & 28.89 & 0.17 & 0.30 & 15485.91 \\
0.50 & ABF-A-D & 27 & 92.07 & 0.11 & 0.26 & 136.18 \\
0.50 & SBF & 0 & 0.00 & 0.45 & 0.45 & - \\
0.50 & SBF-D & 0 & 0.00 & 0.79 & 0.79 & 3030.91 \\
0.50 & SBF-A-D & 41 & 244.02 & 0.01 & 0.15 & 1301.64 \\
1.00 & ABF & 0 & 0.00 & 0.74 & 0.74 & - \\
1.00 & ABF-D & 1 & 370.72 & 0.30 & 0.31 & 20014.76 \\
1.00 & ABF-A-D & 12 & 291.72 & 0.26 & 0.36 & 139.87 \\
1.00 & SBF & 0 & 0.00 & 0.81 & 0.81 & - \\
1.00 & SBF-D & 0 & 0.00 & 1.00 & 1.00 & 2010.27 \\
1.00 & SBF-A-D & 0 & 0.00 & 0.54 & 0.54 & 293.73 \\
\bottomrule
\end{tabular}
\caption{Summary of results for instances grouped by radius values}
\label{tab:general performance r-values}
\end{table}

Figure \ref{fig:profile-r-values} and Table \ref{tab:general performance r-values} present the results grouped by the value of $r$ used in the generation of the instances. 

On instances with small radii, i.e. those generated using $r=0.25$, with any value of $\sigma$, we observe that the three formulations that showed the best performance overall, \ref{abf-a-d}, \ref{sbf-a-d}, and \ref{abf-d}, are also the best performers. However, in this set, \ref{abf-d} and \ref{abf-a-d} perform considerably better than \ref{sbf-a-d}, with \ref{abf-a-d} finding all but 3 optimal solutions, and \ref{abf-d} finding 2 more optimal solutions than \ref{sbf-a-d}, but in considerably lesser times, and with smaller gaps on non-optimal instances. Additionally, we note that \ref{abf} also dominates its sequence-based counterpart \ref{sbf}, which has a performance relatively similar to that of its decomposition \ref{sbf-d}.

On instances with medium radii, i.e. those generated using $r=0.5$, with any value of $\sigma$, we observe a different behavior. In particular, we note that \ref{sbf-a-d} is the formulation that solves the most instances to optimality, but its overall performance is arguably worse than that of \ref{abf-a-d}, which solves only 3 less instances, but has significantly smaller gaps on non-optimal instances. \ref{abf-d} shows a performance slightly worse than that of \ref{abf-a-d}, with similar average gaps, but solving 7 less instances to optimality. We also note that \ref{sbf} seems to perform better than \ref{abf} when it comes to obtaining optimal solutions, getting 4 more than \ref{abf}, but its average gap is larger.

Finally, on instances generated with $r=1$, we observe that \ref{sbf-a-d} presents an exceptional performance, finding optimal solutions for 27 out of 45 instances, while achieving gaps similar to those of \ref{abf-a-d} in the non-optimal cases. In contrast, no other formulation manages to find more than 15 optimal solutions. Once again, \ref{abf-a-d} shows the second best performance on this set. We also observe that, contrary to what happened in the small radii regime, \ref{sbf} has a substantially better performance than \ref{abf}.

Lastly, we note that, in general, the problem seem to become more difficult as the radii increase. This is particularly evident on all three arc-based formulations, which suffer a significant performance drop as the radii increase, while the sequence-based formulations maintain a more consistent performance across different radii, but still showing slightly worse results overall. We theorize that this behavior is due to the fact that, for smaller radii, the problem more closely resembles the classical TSP, where the formulations analog to \ref{abf} are known to perform considerably better than those analog to \ref{sbf}. However, as the radii increase, the problem diverges from the classical TSP, and the advantages of arc-based formulations are less pronounced.

\subsection{Results for approximate formulations}

Since the decompositions based on the approximation of the second-order cone (\ref{sbf-a-d} and \ref{abf-a-d}) are the ones with the best overall performance, it is of interest to analyze the quality of the approximations by themselves, \ref{sbf-approx} and \ref{abf-approx}. To do this, we present in Figure \ref{fig:approx-scatter} scatter plots comparing the runtime and the optimality gap of these two approximate formulations, where each point corresponds to a different instance. The X-axis of the first plot represents the runtime of \ref{abf-approx}, while the Y-axis represents the runtime of \ref{sbf-approx}. The second plot has the same axes, but now they represent the optimality gap of each formulation, i.e. the gap between the lower bound obtained from the approximation, and the lower bound associated to solving the exact SOCP formulation with the visit order fixed according to the solution of the approximation. The diagonal line is the identity; thus, points above this line represent instances where the values associated with \ref{sbf-approx} are greater than those associated with \ref{abf-approx}, while points below the line represent instances where the opposite is true. Different marker shapes and colors are used according to the values of $n$ and $r$ corresponding to each instance. Additionally, Table \ref{tab:approx-performance} summarizes the results of both formulations, where the columns represent the formulation used, the average runtime, the average gap, and the number of instances where the approximation was solved to optimality.

\begin{figure}[h]
\centering
\includegraphics[width=\linewidth]{comparison_scatterplots.png}
\caption{Runtime and gap of approximate formulations}\label{fig:approx-scatter}
\end{figure}

\begin{table}[h]
\centering
\begin{tabular}{lrrr}
\toprule
Formulation & Avg Runtime (s) & Avg Gap (\%) & \# Solved \\
\midrule
ABF-A & 217.51 & 8.61 & 97 \\
SBF-A & 224.97 & 10.74 & 89 \\
\bottomrule
\end{tabular}
\caption{Summary of results for approximate formulations}
\label{tab:approx-performance}
\end{table}


With regards to the runtimes, we observe that \ref{sbf-approx} either solves the instance to optimality in under 200 seconds, or it fails on time. On the other hand, \ref{abf-approx} also solves most instances quickly, but has more variability in this regard. Despite some instances reaching the time limit, both formulations solve most instances in less than 100 seconds, proving the approximations to be tractable in practice.

In terms of optimality gap, we observe that on most instances the gap obtained is the same for both approximations, with values under 2.5\%. In fact, instances where the gaps obtained by the approximations differ are precisely those where one of them fails to find an optimal solution. Lastly, we note that, for both approximations, the harder instances are those with $n=20$, with \ref{abf-approx} having a considerably better performance on instances with small and medium radii out of this set, while the exact opposite happens on instances with large radii.

These results suggest that the approximations are fairly strong, offering high-quality solutions without incurring excessive computational costs, with the arc-based approximation being slightly superior overall to the sequence-based one, which further supports the idea that \ref{abf-a-d} and \ref{sbf-a-d} are the best performing exact formulations due to the quality of the approximations they are based on leading to accurate master problem cost estimations, thus reducing the need for cut generation.


\section{Conclusions}\label{sec8}

This work addresses the Close-Enough Traveling Salesman Problem (CETSP). We present two exact formulations for this problem based on second-order cone programming, along with approximated mixed-integer linear formulations derived from polyhedral approximations of the second-order cone. Additionally, we propose decomposition schemes for both exact and approximate formulations, utilizing cut generation to iteratively refine the master problem. 

Our computational experiments demonstrate the effectiveness of the proposed formulations and decomposition schemes, with the decompositions based on the approximated formulations showing the best overall performance. We also identify different instance regimes where each formulation excels, providing insights into their strengths and weaknesses. Furthermore, we critically analyze the decomposition scheme proposed by \citet{walteros}, highlighting issues with the validity of the cuts used and discussing the implications for convergence to optimality.

Future research could explore the development of stronger cuts for the decomposition schemes, especially considering the enumerative nature of the cuts used in the best-performing decompositions. Additionally, the formulations and decomposition schemes presented could be extended to address related problems, such as the Prize-Collecting TSP with neighborhoods or the Vehicle Routing Problem with neighborhoods. Finally, the decomposition approach based on the approximation of the second-order cone could be adapted to solve other optimization problems involving conic constraints and binary variables, potentially leading to new insights and solution methods in this area.

% =============================================================================

\newpage

\begin{appendices}

\section{Comparison of cuts for \ref{sbf-a-d}}\label{appendix}

In figure \ref{fig:perf-compare-sbf}, we present a performance profile plot comparing the performance of \ref{sbf-a-d} when using cuts \eqref{ll-cut-sbf} and \eqref{ll-cut-sim-sbf}, which we note LL cuts, against the performance when using \eqref{dual-cut-sbf}, which we denote Dual cuts, and the performance when adding both types of cuts. Detailed results are presented in Table \ref{tab:compare-sbf-cuts}.

\begin{figure}[h]
\centering
\includegraphics[width=0.6\linewidth]{sbf-a-d-performance-comparison.png}
\caption{Performance Profile Plot of \ref{sbf-a-d} with different cut types}\label{fig:perf-compare-sbf}
\end{figure}

\begin{table}[h]
\centering
\begin{tabular}{lrrrrr}
\toprule
Cut type & \#Opt & TimeOpt & Gap & Gap (NoOpt) & \#Cuts \\
\midrule
LL cuts & 86 & 119.36 & 0.18 & 0.51 & 599.11 \\
Dual cuts & 81 & 101.49 & 0.19 & 0.47 & 198.27 \\
LL and Dual cuts & 78 & 93.99 & 0.21 & 0.49 & 480.24 \\
\bottomrule
\end{tabular}
\caption{Summary of results for \ref{sbf-a-d} with different cut types}
\label{tab:compare-sbf-cuts}
\end{table}

We observe that there is no performance gain from adding both types of cuts, while the performances when using only one type of cut are fairly similar, with LL cuts having a slight edge overall. It should also be noted that the number of cuts generated when using Dual cuts is considerably smaller than when using LL cuts, which suggests that these cuts are stronger in nature, but the overall performance indicates that the simplicity of LL cuts makes them more effective in practice.

\end{appendices}
% =============================================================================
\newpage

\bibliography{sn-bibliography}% common bib file
%% if required, the content of .bbl file can be included here once bbl is generated
%%\input sn-article.bbl


\end{document}
